\section{The Fire Which Burneth for an Hour}

Among the first Christians to write of hell after the apostles were a bishop on the road to the beasts and an old man before the proconsul's fire. They spoke of the eternal fire because it set every other fear in its place.

When Saint Polycarp, bishop of Smyrna, stood before the proconsul, he would not curse Christ: ``Eighty and six years have I served Him, and He never did me any injury: how then can I blaspheme my King and my Saviour?''\endnote{\emph{The Martyrdom of Polycarp} 9, \emph{The Ante-Nicene Fathers}, vol.~1.} He was threatened first with wild beasts and then with burning, and the Church of Smyrna recorded his answer:

\begin{quote}
Thou threatenest me with fire which burneth for an hour, and after a little is extinguished, but art ignorant of the fire of the coming judgment and of eternal punishment, reserved for the ungodly. But why tarriest thou? Bring forth what thou wilt.\endnote{Ibid., 11.}
\end{quote}

\noindent The same letter says of the martyrs who had gone before him that they ``despised all the torments of this world, redeeming themselves from eternal punishment by [the suffering of] a single hour. For this reason the fire of their savage executioners appeared cool to them. For they kept before their view escape from that fire which is eternal and never shall be quenched.''\endnote{Ibid., 2.}

Saint Ignatius of Antioch, writing to the Ephesians on his way to die at Rome, turned the same fear against the corruption of the faith. ``Do not err, my brethren. Those that corrupt families shall not inherit the kingdom of God. If, then, those who do this as respects the flesh have suffered death, how much more shall this be the case with any one who corrupts by wicked doctrine the faith of God, for which Jesus Christ was crucified! Such an one becoming defiled [in this way], shall go away into everlasting fire, and so shall every one that hearkens unto him.''\endnote{Ignatius of Antioch, \emph{To the Ephesians} 16 (shorter recension), \emph{The Ante-Nicene Fathers}, vol.~1.} The Second Letter of Clement applied Isaias's verse to those who do not keep the seal of their baptism: ``For of those who do not preserve the seal [unbroken], [the Scripture] says, `Their worm shall not die, and their fire shall not be quenched, and they shall be a spectacle to all flesh.'\,''\endnote{Second Letter of Clement 7.}

Already in the second century the pagans mocked the teaching as a scare devised to keep men in line, and Saint Justin answered them. ``And that no one may say what is said by those who are deemed philosophers, that our assertions that the wicked are punished in eternal fire are big words and bugbears, and that we wish men to live virtuously through fear, and not because such a life is good and pleasant; I will briefly reply to this, that if this be not so, God does not exist; or, if He exists, He cares not for men, and neither virtue nor vice is anything.''\endnote{Justin Martyr, \emph{Second Apology} 9, \emph{The Ante-Nicene Fathers}, vol.~1.}

Tertullian made the thought of hell a spur to confession. Writing on repentance, he spoke to the Christian who shrinks from the humiliation of confessing his sins:

\begin{quote}
If you shrink back from exomologesis, consider in your heart the hell, which exomologesis will extinguish for you; and imagine first the magnitude of the penalty, that you may not hesitate about the adoption of the remedy.\endnote{Tertullian, \emph{On Repentance} 12, trans.\ S. Thelwall, \emph{The Ante-Nicene Fathers}, vol.~3, as revised for New Advent, accessed 16 October 2026.}
\end{quote}

\noindent He goes on: ``Why do you desert your own salvation? Why are you tardy to approach what you know heals you?''\endnote{Ibid.} Tertullian asks the penitent to weigh the shame of confession against hell, as Polycarp weighed the proconsul's fire against the eternal one.

The martyrs keep the order of fears that Christ himself gave, when he told his disciples not to fear those who kill the body, but him who can destroy both soul and body in hell. Most of what Christians fear is the fire that burns for an hour: loss of money, of reputation, of comfort, of friends, the mockery of the world, the discomfort of a hard duty. The martyrs feared hell more than all these, and so no threat of the world could move them.
